\documentclass[11pt,a4paper]{article}
\usepackage[margin=1in]{geometry}
\usepackage[T1]{fontenc}
\usepackage[utf8]{inputenc}
\usepackage{lmodern}
\usepackage{microtype}
\usepackage{amsmath,amssymb}
\usepackage{booktabs}
\usepackage{tabularx}
\usepackage{array}
\usepackage{enumitem}
\usepackage[hidelinks]{hyperref}
\usepackage{xcolor}
\setlength{\parindent}{0pt}
\setlength{\parskip}{0.55em}
\renewcommand{\arraystretch}{1.15}
\newcolumntype{Y}{>{\raggedright\arraybackslash}X}
\newcolumntype{L}[1]{>{\raggedright\arraybackslash}p{#1}}
\title{A 3D Medium Operatorial Framework for Environmental Tidal Response:\\Negative Metric-Gate Audit and Operatorial Status Update}
\author{Pietro Risi\\Independent researcher}
\date{Version 3.0.0 status update}
\begin{document}
\maketitle
\begin{abstract}
This version 3.0.0 status update summarizes the post-V1018 correction of the public GTE 3D Medium record. The present version corrects the public metadata, adds a minimal bibliography, and clarifies that the 3D-to-2D relation is an architectural composition of operatorial modules rather than a new bridge result. The scientific content is unchanged. The defensible 3D core remains an operatorial environmental tidal-response pipeline. The negative update is the V993--V1001 audit: with the current assets, no non-circular route from the operatorial medium to a physical metric, universal matter/light coupling, Newtonian limit, or local PPN viability has been derived. The 2D residual is retained only as an operatorial/EFT-inspired environmental memory-screening module. No new numerical execution, fit, detection, gravity derivation, dark-matter claim, or dark-energy claim is made.
\end{abstract}

\section{Relationship to the existing Zenodo version chain}
This document is intended as a corrected version 3.0.0 update for the existing GTE 3D Medium Zenodo version chain, not as a new standalone record. The current public lineage should be preserved through Zenodo's new-version workflow rather than by creating a separate record with duplicated content.

The prior public operatorial core is the V669B/V845 lineage. The version 3.0.0 update should be described as a status and negative-gate update: it adds the V993--V1001 no-go map and corrects the V1018 metadata and bridge wording, while keeping the physical claim boundary closed.

\section{Defensible 3D operatorial core}
The surviving 3D core is the operatorial chain
\[
\rho \rightarrow \Phi \rightarrow E^{\rm proxy}_{ij}/T^{\rm smooth}_{ij} \rightarrow J_{\rm tidal\ grid\ shell} \rightarrow U.
\]
In the earlier public record, the auxiliary response was described by a screened operator
\[
(-\nabla^2 + m_{\rm env}^2)U = J_{\rm env},
\]
with the three-dimensional Green kernel
\[
G_\mu(r) = \frac{e^{-\mu r}}{4\pi r}.
\]
Unless explicitly specified, $U$ denotes an auxiliary operatorial response variable, not a physical gravitational potential. The status of this object remains narrow: it is a reproducible operator-level environmental tidal-response construction. It is not a true relativistic Weyl tensor calculation, not an SMS projected-Weyl calculation, not a physical gravitational field equation, and not a metric derivation.

\section{Negative metric-gate audit: V993--V1001}
The scientific content added after the earlier public record is primarily negative and methodological. V993--V1001 tested several routes for promoting the operatorial medium toward physical gravity. The audited routes included variational/elastic response, intrinsic strain variables, intrinsic connection/curvature, pregeometric-to-metric maps, and bulk/embedding-action routes. The active conclusion is:

\begin{quote}
The framework has not failed as an operatorial environmental-response framework. It has failed, with current assets, as a derivation of physical gravity.
\end{quote}

\begin{tabularx}{\textwidth}{L{0.23\textwidth}L{0.16\textwidth}Y}
\toprule
Gate & Status & Reason \\
\midrule
Physical metric & CLOSED & no non-circular derivation of $g_{\mu\nu}$ \\
Matter/light coupling & CLOSED & no universal coupling to a derived metric \\
Newtonian limit & CLOSED & no target-free derivation of Poisson or $1/r$ behavior \\
Local PPN viability & CLOSED & no $\gamma_{\rm PPN}$ pass \\
True Weyl/SMS & CLOSED & no true relativistic Weyl or 5D projected-Weyl calculation \\
Observational validation & CLOSED & no fit or detection \\
\bottomrule
\end{tabularx}

\section{2D/EFT residual: operatorial module only}
The surviving 2D-inspired residual is only the operatorial/EFT memory-screening schematic
\[
J_{\rm env} := I_L \rightarrow U \rightarrow K_{\rm env} \rightarrow Q_D.
\]
Here again, $U$ denotes an auxiliary response variable unless a separate physical identification is supplied. It is an instance of the same screened-response template, not automatically the same physical field as any variable in the 3D chain. This is not a recovered 2D membrane theory of gravity. The source $I_L$ is a fixed operatorial environmental intensity. The response $U$, kernel $K_{\rm env}$, and diagnostic $Q_D$ remain auxiliary/operatorial unless a separate physical metric map, conservation structure, and matter/light coupling are derived.

\section{Architectural composition, not bridge result}
The V1018 wording treated the 3D-to-2D relationship too prominently. The corrected wording is:
\begin{quote}
The 3D source/intensity output may be used as an input convention for the 2D/EFT-inspired memory-screening module. This is architectural composition between operatorial modules, not a theorem, not a non-trivial compatibility result, and not a physical bridge to gravity.
\end{quote}
Thus the safe composition is
\[
\hbox{3D source/intensity} \rightarrow I_L/J_{\rm env} \rightarrow U/K_{\rm env} \rightarrow Q_D,
\]
but no result attaches to this composition beyond bookkeeping and claim separation.

\section{Claim boundary}
This release claims: a preserved 3D operatorial environmental tidal-response core; a negative V993--V1001 map of failed or blocked routes toward physical gravity; a corrected author/bibliography/wording update for a version 3.0.0 public upload; an operatorial/EFT residual module only at the level of $I_L$, $U$, $K_{\rm env}$, and $Q_D$; and architectural composition between the 3D and 2D/EFT modules, not a scientific bridge result.

This release does not claim: derivation of gravity; derivation of dark matter or dark energy; local PPN viability; a Newtonian or Poisson derivation; a true Weyl/SMS calculation; an observational detection; or a physical replacement for general relativity.

\section*{References}
\begin{enumerate}[leftmargin=*,label={\arabic*.}]
\item P. Risi. \emph{A Membrane/Medium-Inspired Effective Framework for Environmental Tidal Response}. Zenodo, 2026. doi:10.5281/zenodo.20594324.
\item P. Risi. \emph{A 3D Medium Operatorial Framework for Environmental Tidal Response: Complete Working Paper after Z-line All-Sigma and GR Effective-Source Diagnostic Branch}. Zenodo, 2026. doi:10.5281/zenodo.20834009.
\item H. Yukawa. On the interaction of elementary particles. \emph{Proceedings of the Physico-Mathematical Society of Japan}, 17:48--57, 1935.
\item J. Khoury and A. Weltman. Chameleon fields: Awaiting surprises for tests of gravity in space. \emph{Physical Review Letters}, 93:171104, 2004. doi:10.1103/PhysRevLett.93.171104; arXiv:astro-ph/0309300.
\item J. Khoury and A. Weltman. Chameleon cosmology. \emph{Physical Review D}, 69:044026, 2004. doi:10.1103/PhysRevD.69.044026; arXiv:astro-ph/0309411.
\item C. Burrage and J. Sakstein. Tests of chameleon gravity. \emph{Living Reviews in Relativity}, 21:1, 2018. doi:10.1007/s41114-018-0011-x.
\item A. I. Vainshtein. To the problem of nonvanishing gravitation mass. \emph{Physics Letters B}, 39(3):393--394, 1972. doi:10.1016/0370-2693(72)90147-5.
\item O. Hahn, C. Porciani, C. M. Carollo, and A. Dekel. Properties of dark matter haloes in clusters, filaments, sheets and voids. \emph{Monthly Notices of the Royal Astronomical Society}, 375(2):489--499, 2007. doi:10.1111/j.1365-2966.2006.11318.x.
\item J. E. Forero-Romero, Y. Hoffman, S. Gottloeber, A. Klypin, and G. Yepes. A dynamical classification of the cosmic web. \emph{Monthly Notices of the Royal Astronomical Society}, 396(3):1815--1824, 2009. doi:10.1111/j.1365-2966.2009.14885.x.
\end{enumerate}
\end{document}
